% Detailed dual-boot installation procedure.
% Included by main.tex after the pilot's disk plan has been frozen.

\section{Pilot installation --- measure, write, prove}

The profile is a proposed geometry, not permission to erase a disk. Print the
plan, then compare it with the laptop in front of you. Use sectors as the
authoritative unit; rounded GiB figures are only a human cross-check.

\begin{enumerate}
\item Photograph the firmware summary and the physical disk label if it is
accessible. Record firmware revision, disk model, serial and advertised
capacity.
\item Boot the approved PXE release in UEFI mode. Record the release version
and manifest digest shown by the installer.
\item Run the read-only preflight. Save \texttt{lsblk -b -o
NAME,SIZE,MODEL,SERIAL,TYPE,TRAN} and \texttt{wipefs -n} output.
\item Compare the discovered logical-sector size, last usable sector and disk
serial with the frozen plan. Differences are a stop condition.
\end{enumerate}

\askbefore{Is this the one authorized internal disk? Is the PXE or USB medium
absent from the target list? Does the disk have at least the Windows and Arch
minimums after the ESP, Microsoft reserved partition and recovery allowance?
Are there files on the existing disk that have not been independently
recovered and opened?}

\begin{tabularx}{\linewidth}{@{}p{.24\linewidth}p{.22\linewidth}p{.24\linewidth}L@{}}
\toprule
\textbf{Measurement} & \textbf{Frozen value} & \textbf{Observed value} & \textbf{Pass rule}\\
\midrule
Disk model and serial & \placeholder{plan} & \placeholder{preflight} & Exact serial; model consistent\\
Disk bytes & \placeholder{plan} & \placeholder{lsblk} & Exact byte count\\
Logical sector bytes & \placeholder{plan} & \placeholder{blockdev} & Exact match\\
First/last usable LBA & \placeholder{plan} & \placeholder{sgdisk} & Exact match\\
PXE release digest & \placeholder{release} & \placeholder{screen} & Exact SHA-256\\
\bottomrule
\end{tabularx}

\begin{gate}{Last reversible point}
Do not type the disk serial until the table passes. The next operation destroys
the selected partition table. No hostname, device name such as
\texttt{/dev/nvme0n1}, or general confirmation substitutes for the complete
serial.
\end{gate}

\subsection{Create and verify the GPT}

The automation creates, in order, one shared EFI System Partition (ESP), one
Microsoft Reserved partition (MSR), the Windows partition, Windows recovery,
and the Arch root partition. Start each resizable partition on the alignment
boundary emitted by the profile. Do not calculate boundaries by eye in a GUI.

\begin{center}
\begin{tikzpicture}[x=1cm,y=1cm,font=\scriptsize]
  \draw[wf pencil] (-5.75,-.20) rectangle (5.70,1.10);
  \draw[wf faint] (-5.22,-.20)--(-5.22,1.10);
  \draw[wf faint] (-4.93,-.20)--(-4.93,1.10);
  \draw[wf faint] (1.25,-.20)--(1.25,1.10);
  \draw[wf faint] (2.05,-.20)--(2.05,1.10);
  \node[rotate=90] at (-5.49,.45) {ESP};
  \node[rotate=90] at (-5.08,.45) {MSR};
  \node at (-1.84,.45) {Windows};
  \node[rotate=90] at (1.65,.45) {WinRE};
  \node at (3.87,.45) {Arch root};
  \draw[wf pencil,<->] (-5.75,-.47)--(-5.22,-.47);
  \node[below] at (-5.49,-.48) {1 GiB};
  \draw[wf pencil,<->] (-4.93,-.47)--(1.25,-.47);
  \node[below] at (-1.84,-.48) {profile-derived};
  \draw[wf pencil,<->] (2.05,-.47)--(5.70,-.47);
  \node[below] at (3.87,-.48) {profile-derived};
\end{tikzpicture}
\wfcaption{One GPT and one shared ESP; exact LBA boundaries come from the signed plan.}
\end{center}

\expectthis{A second read of the GPT reports the planned partition count,
type GUIDs, start sectors and end sectors. No filesystem signature exists
outside the partitions declared by the profile.}
\failurebranch{If any start sector, type GUID, serial or disk byte count differs,
stop before formatting. Save both observations, remove installer power, and
regenerate the plan for the hardware actually present.}

\subsection{Install Windows first}

Boot Windows Setup from the versioned PXE entry. Load only the driver bundle
whose hardware identifiers and digest were approved with the release. At the
disk chooser, select the pre-created partition labelled for Windows; do not
delete the table and do not allow Setup to consume the Arch partition.

\begin{enumerate}
\item Confirm Setup reports the planned Windows partition size within its
display-rounding tolerance. Ask: ``Which partition will Setup format, and what
two observed facts identify it?''
\item Install Windows 11 Pro. Firmware activation may occur later; activation
status is not an installation gate.
\item At the first desktop, install firmware and drivers from the laptop
vendor or Windows Update. Record each device that retains a warning icon.
\item Apply all update rounds, restarting and checking again until two
successive scans report no required update. Optional preview drivers remain
excluded unless a recorded hardware fault requires one.
\end{enumerate}

\proofpair{Get-Disk | Format-Table Number,FriendlyName,SerialNumber,Size}{The
system disk serial and byte count identify the authorized disk.}
\proofpair{Get-Partition | Sort-Object Offset}{Partition offsets and sizes
match the frozen plan; no unexpected data partition exists.}
\proofpair{Get-PnpDevice | Where-Object Status -ne 'OK'}{Returns no unexplained
device; record any intentional disabled device.}
\proofpair{Confirm-SecureBootUEFI}{Reports \texttt{False}, consistent with the
phase-1 boundary; an unsupported result is investigated, not assumed.}
\proofpair{manage-bde -status}{All workstation volumes report protection off;
device encryption must not have silently enabled itself.}

\begin{gate}{Windows recovery proof}
Open the Windows Recovery Environment once before installing Arch. Confirm the
keyboard works and that Startup Repair and Command Prompt are visible. Exit
without changing the disk. This proves the recovery partition is bootable
before another operating system changes UEFI entries.
\end{gate}

\subsection{Install Arch into its declared partition}

Return to the same PXE release and select Arch. Re-run disk identity preflight;
a successful Windows boot does not waive serial authorization. Format only the
Arch root partition. Mount the existing ESP at \texttt{/efi}; never format it.

\begin{enumerate}
\item Check time synchronization before package installation. Record clock
offset and time source; Kerberos will depend on this later.
\item Mount Arch root, then mount the existing FAT32 ESP. List
\texttt{/efi/EFI/Microsoft/Boot/bootmgfw.efi}; its presence is a stop-check
before installing a boot manager.
\item Install the base system, Intel microcode, NetworkManager, an LTS-capable
kernel path, and the approved graphics, Wi-Fi and firmware packages.
\item Generate \texttt{/etc/fstab} by UUID. Inspect every line: root must use
the Arch filesystem UUID and the ESP must use the shared FAT UUID.
\item Install \texttt{systemd-boot} without removing the Microsoft directory.
Create the Arch entry, set a five-second timeout, and set Windows
\texttt{auto-windows} as the default.
\end{enumerate}

\proofpair{findmnt --verify --verbose}{Reports no fstab parse or unreachable
source error.}
\proofpair{bootctl status}{Shows the shared ESP, an Arch entry and the automatic
Windows entry; default is Windows and timeout is five seconds.}
\proofpair{Verify the Windows EFI loader exists}{Confirm
\texttt{/efi/EFI/Microsoft/Boot/bootmgfw.efi} remains present after
the Arch boot manager is installed.}
\proofpair{lspci -k}{Each storage, graphics, Ethernet and wireless controller
has the expected kernel driver bound.}
\proofpair{systemctl --failed}{Returns no unexplained failed unit after a native
Arch boot.}

\failurebranch{If the automatic Windows entry is absent, do not hand-edit or
copy Microsoft binaries. Verify the ESP mount and the original Microsoft path.
If Windows files are missing, use Windows recovery media to repair them, then
rerun \texttt{bootctl status}.}

\subsection{Prove both native boot paths}

Power the laptop fully off between tests. A reboot alone may preserve firmware
or device state and conceal a cold-start defect.

\begin{enumerate}
\item Power on without touching a key. Measure from menu appearance until
Windows begins loading. Pass only if Windows is selected after five seconds.
\item Restart, select Arch, and measure time to a local login prompt. Record
network association time separately from operating-system boot time.
\item Use the firmware's one-time boot menu to start Windows Boot Manager
directly, then to start Linux Boot Manager directly.
\item Repeat once with Ethernet disconnected. Both local operating systems must
reach login independently of PXE, AD, DNS and network storage.
\end{enumerate}

\begin{tabularx}{\linewidth}{@{}p{.27\linewidth}p{.20\linewidth}p{.24\linewidth}L@{}}
\toprule
\textbf{Boot path} & \textbf{Expected} & \textbf{Measured} & \textbf{Evidence}\\
\midrule
Default menu & Windows, 5 s & & video / stopwatch\\
Selected Arch & local login & & journal boot ID\\
Firmware Windows entry & Windows login & & firmware photo\\
Firmware Linux entry & Arch login & & firmware photo\\
Both, cable absent & local login & & two boot IDs\\
\bottomrule
\end{tabularx}

\subsection{Recovery order}

Recovery preserves the other operating system before it improves convenience.
Do not repartition to repair a boot menu.

\begin{enumerate}
\item \textbf{Firmware first:} open the one-time menu. If either native entry
works, boot it and record \texttt{efibootmgr -v} or
\texttt{bcdedit /enum firmware}.
\item \textbf{Selector second:} from Arch media, mount the existing root and
ESP, confirm their UUIDs, chroot, then reinstall systemd-boot and regenerate
the Arch entry. Do not format either filesystem.
\item \textbf{Windows files third:} if the Microsoft EFI binary is absent or
Windows reports boot damage, use the matching Windows 11 recovery image and
\texttt{bcdboot} against the existing Windows and ESP volumes.
\item \textbf{Rollback last:} select the previous verified PXE release if a
current installer or driver bundle is suspect. A PXE rollback does not modify
an already installed workstation.
\end{enumerate}

\askbefore{Which layer failed: firmware entry, selector file, OS loader,
filesystem, or operating system? Which read-only observation proves that
classification? Can the repair be performed without formatting or moving a
partition? Where is the recovery evidence saved?}

\begin{gate}{Escalation boundary}
If neither filesystem is readable, the disk reports new media errors, or the
observed partition boundaries differ from the signed plan, stop. Image the disk
or replace the disposable pilot; do not use filesystem repair, partition
reconstruction, or installation retry as an exploratory diagnostic.
\end{gate}
